%% Question bank chapter file - generated by bank_tools.py build (do not edit while the bank is still being built:
%% it is regenerated from the per-test data in Question Bank\data). Every question is one bankq block:
%%   \begin{bankq}[id=..., source={...}, paper=.., qno=.., type=..., subtopic={...}, answer={...}, answer_src=key|solved,
%%                 working={...}, difficulty=1-10, diff_parts={K S I C F}, diff_note={...},
%%                 group=<passage id or empty>, also_in={duplicate copies in other tests}]
%% difficulty = min(10, K+S+I+C+F-1): K concepts 1-3, S steps 1-4, I insight 0-2, C computation 0-2, F format 0-1
%%   ...question LaTeX, options, figures (\bfig / \bfigi -> img/)...
%%   \end{bankq}
%% Passages shared by several questions are bankpassage blocks; their questions carry group=<passage id>.
\documentclass[10pt]{article}
\usepackage[a4paper,margin=1.8cm]{geometry}
\usepackage[T1]{fontenc}
\usepackage{amsmath,amssymb,mathtools}
\usepackage{graphicx}
\usepackage[export]{adjustbox}
\usepackage{xcolor,needspace,keyval}
\usepackage{enumitem}
\usepackage{booktabs,array,tabularx}
\usepackage[hidelinks]{hyperref}
\setlength{\parindent}{0pt}
\setlength{\parskip}{3pt}
\makeatletter
\define@key{bq}{id}{\def\bq@id{#1}}\define@key{bq}{source}{\def\bq@src{#1}}\define@key{bq}{paper}{\def\bq@paper{#1}}
\define@key{bq}{qno}{\def\bq@qno{#1}}\define@key{bq}{type}{\def\bq@type{#1}}\define@key{bq}{subtopic}{\def\bq@sub{#1}}
\define@key{bq}{answer}{\def\bq@ans{#1}}\define@key{bq}{answer_src}{\def\bq@asrc{#1}}\define@key{bq}{working}{\def\bq@work{#1}}
\define@key{bq}{group}{\def\bq@group{#1}}\define@key{bq}{also_in}{\def\bq@also{#1}}
\define@key{bq}{difficulty}{\def\bq@diff{#1}}\define@key{bq}{diff_parts}{\def\bq@dparts{#1}}\define@key{bq}{diff_note}{\def\bq@dnote{#1}}
\newenvironment{bankq}[1][]{\def\bq@work{}\def\bq@also{}\def\bq@group{}\def\bq@diff{?}\def\bq@dparts{}\def\bq@dnote{}\setkeys{bq}{#1}%
  \par\needspace{4\baselineskip}\medskip\noindent{\small\textbf{\bq@id}\hfill\color{gray}[\bq@src, Paper \bq@paper, Q\bq@qno]
  \bq@type\ $\cdot$ \bq@sub\ $\cdot$ \textbf{D\bq@diff}}\par\nopagebreak\relax}
  {\par{\small\color{teal}Answer: \textbf{\bq@ans}\ifx\bq@work\@empty\else\ (solved: \bq@work)\fi
  \ifx\bq@also\@empty\else\par\color{gray}Also in: \bq@also\fi}\par}
\define@key{bp}{id}{\def\bp@id{#1}}\define@key{bp}{source}{\def\bp@src{#1}}\define@key{bp}{questions}{\def\bp@qs{#1}}
\newenvironment{bankpassage}[1][]{\setkeys{bp}{#1}\par\needspace{4\baselineskip}\bigskip\noindent
  {\small\textbf{Passage \bp@id}\hfill\color{gray}[\bp@src; questions \bp@qs]}\par\nopagebreak\relax\itshape}{\par}
\makeatother
\newcommand{\bfig}[2]{\par\smallskip{\centering\includegraphics[width=#1pt,max width=0.62\linewidth,max height=0.32\textheight]{img/#2}\par}\smallskip}
\newcommand{\bfigi}[2]{\includegraphics[width=#1pt,max width=\linewidth,max height=0.32\textheight,valign=c]{img/#2}}
\newcommand{\bmissing}[1]{\textit{[figure missing in the source conversion: \detokenize{#1}]}}
\newenvironment{qoptions}{\begin{itemize}[leftmargin=1.8em,itemsep=1pt,topsep=2pt,label={}]}{\end{itemize}}
\newcommand{\opt}[2]{\item[(#1)] #2}
\providecommand{\fwstart}{}\providecommand{\fwend}{}\providecommand{\qfullsplit}{}
\newcolumntype{L}[1]{>{\raggedright\arraybackslash}p{#1}}
\newcolumntype{C}[1]{>{\centering\arraybackslash}p{#1}}
\newcolumntype{R}[1]{>{\raggedleft\arraybackslash}p{#1}}
\providecommand{\sectionhdr}[1]{}\providecommand{\subsectionhdr}[1]{}


\begin{document}
{\Large\bfseries Alcohol, Ether, Phenol --- Question Bank}\par\smallskip
{\small Chemistry chapter C19. Questions from all converted test papers, grouped by source series. Answers: printed key where it matches the paper, otherwise solved (answer\_src=solved, with working). Exact repeats across papers appear once, with also\_in listing the other copies.}\par\medskip\hrule

\section*{Score-I papers}
\begin{bankq}[id=SI12-P1-CHEM-S1-Q04, source={Score-I - Test 12 (09 Mar 2022)}, paper=1, qno=4, type=multi, subtopic={Properties}, answer={A,B,C}, answer_src=key, difficulty=4, diff_parts={K2 S2 I0 C0 F1}, diff_note={concepts: acidity and dehydration of alcohols, carbocation stability | steps: acidic strength order; dehydration rate order; major product of dehydration; Lucas reaction rate}]
\bfig{150.6}{SI12_p1_chem_s1_q04_1.png}
Correct option is/are:
\begin{qoptions}
\opt{A}{H$_2$O $>$ P $>$ R \ (Acidic strength)}
\opt{B}{Q $>$ P \ (Rate of dehydration)}
\opt{C}{All P, Q \& R gives same dehydration major product with conc. H$_2$SO$_4$}
\opt{D}{All P, Q \& R react with Lucas reagent with same rate}
\end{qoptions}
\end{bankq}

\section*{TOAS papers}
\begin{bankpassage}[id=TOAS01-P1-CHEM-S3-G5-6, source={TOAS - Test 01 (30 Jan 2022)}, questions={5-6}]
Dehydration of alcohol is an elimination reaction, which may proceed via formation of carbocation intermediate \& carbocation rearrangement is also considered if required. Alkene is the product of the reaction where usually more stable alkene is major product.
\bfig{150.6}{TOAS01_p1_chem_s3_para5_6_1.png}
\end{bankpassage}
\begin{bankq}[id=TOAS01-P1-CHEM-S3-Q05, source={TOAS - Test 01 (30 Jan 2022)}, paper=1, qno=5, type=numerical, subtopic={Properties}, answer={7}, answer_src=key, difficulty=2, diff_parts={K1 S1 I0 C0 F1}, diff_note={concepts: alcohol reactions and alpha hydrogens | steps: major product of the reaction; count alpha hydrogens}, group=TOAS01-P1-CHEM-S3-G5-6]
Number of $\alpha$-H present in major product :
\end{bankq}
\begin{bankq}[id=TOAS01-P1-CHEM-S3-Q06, source={TOAS - Test 01 (30 Jan 2022)}, paper=1, qno=6, type=numerical, subtopic={Properties}, answer={5}, answer_src=key, difficulty=2, diff_parts={K1 S1 I0 C0 F1}, diff_note={concepts: reaction mechanism steps | steps: mechanism of the reaction; count transition states in each step}, group=TOAS01-P1-CHEM-S3-G5-6]
Number of T.S. (transition state) in given reaction for major product :
\end{bankq}
\begin{bankq}[id=TOAS02-P1-CHEM-S2-Q09, source={TOAS - Test 02 (13 Feb 2022)}, paper=1, qno=9, type=multi, subtopic={Properties}, answer={A,B,D}, answer_src=solved, working={SOCl$_2$ alone: retention (S$_N$i); SOCl$_2$/pyridine and PCl$_3$/pyridine: inversion (S$_N$2); HI on 2,2-dimethylcyclohexanol gives a rearranged mixture (C wrong)}, difficulty=3, diff_parts={K1 S2 I0 C0 F1}, diff_note={concepts: substitution of alcohols with SOCl2 and PCl3 and HI | steps: SOCl2 alone and with pyridine; PCl3 with pyridine; HI with rearrangement}]
Which of the following reaction is showing correct major product :
\begin{qoptions}
\opt{A}{\bfigi{235.3}{TOAS02_p1_chem_s2_q09_A.png}}
\opt{B}{\bfigi{244.2}{TOAS02_p1_chem_s2_q09_B.png}}
\opt{C}{\bfigi{168.9}{TOAS02_p1_chem_s2_q09_C.png}}
\opt{D}{\bfigi{245.2}{TOAS02_p1_chem_s2_q09_D.png}}
\end{qoptions}
\end{bankq}
\begin{bankq}[id=TOAS02-P1-CHEM-S3-Q03, source={TOAS - Test 02 (13 Feb 2022)}, paper=1, qno=3, type=numerical, subtopic={Properties}, answer={3}, answer_src=solved, working={Immediate turbidity (Lucas) for 3$^\circ$ and benzylic alcohols: (b), (c), (d); (a) 1$^\circ$, (e) bridgehead, (f) enol, (g) antiaromatic cation do not react}, difficulty=3, diff_parts={K1 S2 I0 C0 F1}, diff_note={concepts: Lucas test | steps: classify each alcohol as primary, secondary, tertiary or benzylic; carbocation stability for bridgehead and enol; count}]
How many of the following compound will form immediate turbidity with ($\mathrm{HCl/ZnCl_2}$):
\bfig{185.2}{TOAS02_p1_chem_s3_q03_1.png}
\end{bankq}
\begin{bankq}[id=TOAS04-P1-CHEM-S2-Q09, source={TOAS - Test 04 (15 Mar 2022)}, paper=1, qno=9, type=multi, subtopic={Properties}, answer={A,B,C}, answer_src=key, difficulty=3, diff_parts={K1 S2 I0 C0 F1}, diff_note={concepts: reagents for conversion of alcohols | steps: PCl5; SOCl2; PCl3; HCl/ZnCl2}]
\bfigi{209.1}{TOAS04_p1_chem_s2_q09_1.png}

Correct pair of $X$ \& $Y$ respectively is :
\begin{qoptions}
\opt{A}{$X=\mathrm{PCl_5}$, $Y=$ \bfigi{88.0}{TOAS04_p1_chem_s2_q09_2.png}}
\opt{B}{$X=\mathrm{SOCl_2}$, $Y=$ \bfigi{88.2}{TOAS04_p1_chem_s2_q09_3.png}}
\opt{C}{$X=\mathrm{PCl_3}$, $Y=$ \bfigi{85.8}{TOAS04_p1_chem_s2_q09_4.png}}
\opt{D}{$X=\mathrm{HCl/ZnCl_2}$, $Y=$ \bfigi{84.2}{TOAS04_p1_chem_s2_q09_5.png}}
\end{qoptions}
\end{bankq}
\begin{bankq}[id=TOAS05-P1-CHEM-S3-Q01, source={TOAS - Test 05 (03 Apr 2022)}, paper=1, qno=1, type=numerical, subtopic={Preparation}, answer={4.00}, answer_src=key, difficulty=3, diff_parts={K1 S2 I0 C0 F1}, diff_note={concepts: Williamson ether synthesis | steps: tertiary or aryl halides cannot react; check each compound; count}]
Number of following compounds which can't be synthesised in good yield by Williamson's reaction
\bfig{293.1}{TOAS05_p1_chem_s3_q01_1.png}
\end{bankq}
\begin{bankq}[id=TOAS05-P2-CHEM-S1-Q02, source={TOAS - Test 05 (03 Apr 2022)}, paper=2, qno=2, type=single, subtopic={Properties}, answer={C}, answer_src=key, difficulty=1, diff_parts={K1 S1 I0 C0 F0}, diff_note={concepts: major product of a reaction | steps: identify the reaction type; major product}]
The major product formed in the following reaction is :
\bfig{81.1}{TOAS05_p2_chem_s1_q02_1.png}
\begin{qoptions}
\opt{A}{\bfigi{45.5}{TOAS05_p2_chem_s1_q02_A.png}}
\opt{B}{\bfigi{58.3}{TOAS05_p2_chem_s1_q02_B.png}}
\opt{C}{\bfigi{46.3}{TOAS05_p2_chem_s1_q02_C.png}}
\opt{D}{\bfigi{53.4}{TOAS05_p2_chem_s1_q02_D.png}}
\end{qoptions}
\end{bankq}
\begin{bankq}[id=TOAS05-P2-CHEM-S1-Q06, source={TOAS - Test 05 (03 Apr 2022)}, paper=2, qno=6, type=single, subtopic={Properties}, answer={B}, answer_src=key, difficulty=3, diff_parts={K1 S2 I1 C0 F0}, diff_note={concepts: cleavage of glycerol with HI | steps: substitution of each OH; reduction of iodide product; count HI | I: glycerol reduces to propene by HI}]
How many moles of HI will be consumed per mole of glycerol in the following reaction.
\bfig{75.6}{TOAS05_p2_chem_s1_q06_1.png}
\begin{qoptions}
\opt{A}{4}
\opt{B}{5}
\opt{C}{6}
\opt{D}{7}
\end{qoptions}
\end{bankq}
\begin{bankq}[id=TOAS05-P2-CHEM-S2-Q01, source={TOAS - Test 05 (03 Apr 2022)}, paper=2, qno=1, type=numerical, subtopic={Properties}, answer={5}, answer_src=solved, working={Immediate turbidity (3$^\circ$ or highly stabilised cation): cyclopropylmethanol, cyclopropenol (aromatic cyclopropenyl cation), tert-butanol, benzyl alcohol, ethoxymethanol; not cyclopentadienol (antiaromatic cation) or cyclohexylmethanol}, difficulty=3, diff_parts={K1 S2 I0 C0 F1}, diff_note={concepts: Lucas test | steps: classify each alcohol; carbocation stability including aromatic cations; count}]
\bfig{123.4}{TOAS05_p2_chem_s2_q01_1.png}
Identify total number of compound gives immidiate turbidity with Lucas reagent.
\end{bankq}
\begin{bankq}[id=TOAS08-P1-CHEM-S3-Q07, source={TOAS - Test 08 (15 May 2022)}, paper=1, qno=7, type=numerical, subtopic={Properties}, answer={7}, answer_src=solved, working={Propan-2-ol to acetone: Cu/300 C, Jones, KMnO$_4$/OH$^-$, Oppenauer (Al(OCMe$_3$)$_3$/quinone), Collins, TsCl then DMSO (Kornblum), PCC; O$_3$ and Zn-Hg/HCl cannot: 7}, difficulty=3, diff_parts={K1 S2 I0 C0 F1}, diff_note={concepts: oxidation of a secondary alcohol | steps: oxidising reagents; reagents that reduce or cleave; count suitable ones}]
How many of reagents can carry out following transformation :
\bfig{148.7}{TOAS08_p1_chem_s3_q07_1.png}
(a) Cu, 300$^\circ$C\newline
(b) $\mathrm{O_3}$ / Zn, $\mathrm{H_2O}$\newline
(c) Zn--Hg / HCl\newline
(d) Jones reagent $\mathrm{CrO_3+H_2O}$\newline
(e) $\mathrm{KMnO_4}$, NaOH, $\Delta$\newline
\bfigi{115.3}{TOAS08_p1_chem_s3_q07_2.png}\newline
(g) Collin's reagent\newline
(h) (i) TsCl, (ii) DMSO, $\mathrm{NaHCO_3}$ , $\Delta$\newline
(i) $\mathrm{C_5H_5HCrO_3Cl^-}$ (PCC)
\end{bankq}
\begin{bankq}[id=TOAS09-P2-CHEM-S1-Q11, source={TOAS - Test 09 (30 May 2022)}, paper=2, qno=11, type=single, subtopic={Properties}, answer={A}, answer_src=key, difficulty=1, diff_parts={K1 S1 I0 C0 F0}, diff_note={concepts: periodate cleavage | steps: vicinal diols are cleaved; the option without adjacent OH groups}]
Which of the following will not undergo periodic cleavage with HIO$_4$, $\Delta$?
\begin{qoptions}
\opt{A}{\bfigi{84.6}{TOAS09_p2_chem_s1_q11_A.png}}
\opt{B}{\bfigi{48.3}{TOAS09_p2_chem_s1_q11_B.png}}
\opt{C}{\bfigi{49.5}{TOAS09_p2_chem_s1_q11_C.png}}
\opt{D}{\bfigi{52.0}{TOAS09_p2_chem_s1_q11_D.png}}
\end{qoptions}
\end{bankq}
\begin{bankq}[id=TOAS09-P2-CHEM-S1-Q13, source={TOAS - Test 09 (30 May 2022)}, paper=2, qno=13, type=single, subtopic={Preparation}, answer={B}, answer_src=key, difficulty=1, diff_parts={K1 S1 I0 C0 F0}, diff_note={concepts: Williamson ether synthesis | steps: alkyl halide must be primary and aryl halides cannot be used; test each ether}]
Which of the following ether can not be synthesised by Williamson ether synthesis in good yield :
\begin{qoptions}
\opt{A}{Ph--O--CH$_2$--CH$_3$}
\opt{B}{\bfigi{69.3}{TOAS09_p2_chem_s1_q13_B.png}}
\opt{C}{CH$_3$--O--CH$_3$}
\opt{D}{Ph--CH$_2$--O--Ph}
\end{qoptions}
\end{bankq}
\begin{bankq}[id=TOAS10-P1-CHEM-S1-Q07, source={TOAS - Test 10 (12 Jun 2022)}, paper=1, qno=7, type=multi, subtopic={Properties}, answer={A,C}, answer_src=key, difficulty=3, diff_parts={K1 S2 I0 C0 F1}, diff_note={concepts: reactions of a diol sequence | steps: identify A from the reagents; oxidation and cleavage products; Tollens and iodoform tests}]
\bfig{221.8}{TOAS10_p1_chem_s1_q07_1.png}
Correct statement(s) is/are
\begin{qoptions}
\opt{A}{(A) is vicinol diol}
\opt{B}{(B) is a dicarboxylic acid}
\opt{C}{(B) gives Tollen's test}
\opt{D}{(D) gives Iodoform test}
\end{qoptions}
\end{bankq}
\begin{bankq}[id=TOAS11-P1-CHEM-S2-Q07, source={TOAS - Test 11 (27 Jun 2022)}, paper=1, qno=7, type=multi, subtopic={Properties}, answer={B,D}, answer_src=key, difficulty=3, diff_parts={K1 S2 I0 C0 F1}, diff_note={concepts: reactions of alcohols and Grignard products | steps: identify P, Q, R, S, T from the figure; Lucas test; isomerism of products}]
For the following reaction the correct option is/are :
\bfig{245.2}{TOAS11_p1_chem_s2_q07_1.png}
\begin{qoptions}
\opt{A}{Both (R) \& (T) gives immediate turbidity with HCl + anhy ZnCl$_2$}
\opt{B}{(P) shows optical isomerism}
\opt{C}{(S) shows geometrical isomerism}
\opt{D}{(Q) is known as grignard reagent}
\end{qoptions}
\end{bankq}
\begin{bankq}[id=TOAS11-P2-CHEM-S1-Q04, source={TOAS - Test 11 (27 Jun 2022)}, paper=2, qno=4, type=multi, subtopic={Preparation}, answer={A,B,C,D}, answer_src=key, difficulty=3, diff_parts={K1 S2 I0 C0 F1}, diff_note={concepts: primary alcohol formation | steps: Grignard with HCHO; hydrogenation of aldehyde; acid chloride with NaBH4; hydroboration}]
Which of the following give 1$^\circ$ alcohol as a major product :
\begin{qoptions}
\opt{A}{H--CHO $+$ CH$_3$MgBr $\to$ $\xrightarrow{\mathrm{HOH/H^{\oplus}}}$}
\opt{B}{CH$_3$--CH$_2$--CHO $\xrightarrow{\mathrm{H_2/Pd}}$}
\opt{C}{CH$_3$--COCl $\xrightarrow[\mathrm{H_2O}]{\mathrm{NaBH_4}}$}
\opt{D}{CH$_3$--CH=CH$_2$ $\xrightarrow[\text{(ii) NaOH / H}_2\text{O}_2]{\text{(i) B}_2\text{H}_6\text{ / THF}}$}
\end{qoptions}
\end{bankq}
\begin{bankq}[id=TOAS12-P1-CHEM-S2-Q07, source={TOAS - Test 12 (10 Jul 2022)}, paper=1, qno=7, type=numerical, subtopic={Properties}, answer={2.00}, answer_src=key, difficulty=2, diff_parts={K1 S1 I0 C0 F1}, diff_note={concepts: reaction with HI | steps: identify the groups in the figure; count the moles of HI per mole}]
\bfig{113.7}{TOAS12_p1_chem_s2_q07_1.png}
How many ($x$) moles of HI consumed in the given reaction ?
\end{bankq}
\begin{bankpassage}[id=TOAS14-P1-CHEM-S2-G5-6, source={TOAS - Test 14 (07 Aug 2022)}, questions={5-6}]
In a polyhydric alcohol, the total number of hydroxyl groups are estimated by the following methods:\par
(i) A known mass of the acetyl derivative of polyhdric alcohol is taken. It is mixed with an excess volume of a standard alkali solution. Both are refluxed. Alkali is used in the hydrolysis of acetyl derivative. The unused alkali is measured by doing titration with a standard acid solution and total number of hydroxyl goups (n) can thus be calculated stoichiometrically.\par
{\footnotesize\noindent A(OH)$_n$ $\xrightarrow{n\text{CH}_3\text{COCl}}$ A(OCOCH$_3$)$_n$ $\xrightarrow{n\text{KOH}}$ A(OH)$_n$ + $n$CH$_3$COOK\par}
{\footnotesize\noindent Polyhydric alcohol \hfill Acetyl derivative\par}
(ii) Total number of-OH groups in vicinal position can be calculated by periodic acid or lead tetra acetate oxidation.
\bfig{141.1}{TOAS14_p1_chem_s2_para5_6_2.png}
\end{bankpassage}
\begin{bankq}[id=TOAS14-P1-CHEM-S2-Q05, source={TOAS - Test 14 (07 Aug 2022)}, paper=1, qno=5, type=numerical, subtopic={Properties}, answer={3.00}, answer_src=key, difficulty=4, diff_parts={K1 S2 I0 C1 F1}, diff_note={concepts: saponification stoichiometry of an acetylated alcohol | steps: mass per mole of the derivative is 92 + 42 n; moles of KOH equal n times moles of derivative; solve n}, group=TOAS14-P1-CHEM-S2-G5-6]
0.436 g of acetyl derivative of a polyhydric alcohol (molecular mass = 92) requires 0.336 g KOH for hydrolysis. The total numbers of hydroxyl group(s) in the alcohol are
\end{bankq}
\begin{bankq}[id=TOAS14-P1-CHEM-S2-Q06, source={TOAS - Test 14 (07 Aug 2022)}, paper=1, qno=6, type=numerical, subtopic={Properties}, answer={2.00}, answer_src=key, difficulty=2, diff_parts={K1 S1 I0 C0 F1}, diff_note={concepts: acetylation mass change | steps: each acetyl group adds 42 to the molar mass 106 so two are needed}, group=TOAS14-P1-CHEM-S2-G5-6]
A compound with molecular formula C$_4$H$_{10}$O$_3$ is converted by the action of acetyl chloride to a compound(s) with molecular weight 190. How many number of acetyl group present in compound (B)
\end{bankq}
\begin{bankq}[id=TOAS16-P2-CHEM-S1-Q06, source={TOAS - Test 16 (04 Sep 2022)}, paper=2, qno=6, type=multi, subtopic={Properties}, answer={C,D}, answer_src=key, difficulty=5, diff_parts={K2 S2 I1 C0 F1}, diff_note={concepts: Williamson-type substitution and acylation, reduction | steps: phenoxide formation; reaction with chloroacetyl chloride; second ring closure; reduction}]
\bfig{236.8}{TOAS16_p2_chem_s1_q06_1.png}
Select correct statement :
\begin{qoptions}
\opt{A}{Q is \bfigi{59.9}{TOAS16_p2_chem_s1_q06_2.png}}
\opt{B}{R is \bfigi{59.9}{TOAS16_p2_chem_s1_q06_3.png}}
\opt{C}{S is \bfigi{117.0}{TOAS16_p2_chem_s1_q06_4.png}}
\opt{D}{Q is \bfigi{79.9}{TOAS16_p2_chem_s1_q06_5.png}}
\end{qoptions}
\end{bankq}
\begin{bankq}[id=TOAS17-P1-CHEM-S1-Q01, source={TOAS - Test 17 (18 Sep 2022)}, paper=1, qno=1, type=multi, subtopic={Properties}, answer={A,B,C}, answer_src=key, difficulty=4, diff_parts={K2 S2 I0 C0 F1}, diff_note={concepts: reduction, oxidation of alcohols | steps: identify the reagents in sequence; structure of the final product; apply each test}]
Which of the following is/are correct for the final product of the given sequence of reaction.
\bfig{245.2}{TOAS17_p1_chem_s1_q01_1.png}
\begin{qoptions}
\opt{A}{Compound P$_3$ on oxidation with (CrO$_3$ + H$_2$SO$_4$) gives a compound which can give positive 2,4-DNP test.}
\opt{B}{Compound P$_3$ on reaction with I$_2$ + NaOH gives yellow precipitate.}
\opt{C}{Compound P$_3$ can react with PCC.}
\opt{D}{Compound P$_3$ on reaction with MnO$_2$ gives carboxylic acid.}
\end{qoptions}
\end{bankq}
\begin{bankq}[id=TOAS17-P1-CHEM-S1-Q02, source={TOAS - Test 17 (18 Sep 2022)}, paper=1, qno=2, type=multi, subtopic={Preparation}, answer={A,B,C,D}, answer_src=key, difficulty=5, diff_parts={K2 S2 I1 C0 F1}, diff_note={concepts: nucleophilic addition of Grignard, intramolecular substitution | steps: Grignard adds to the carbonyl; alkoxide displaces the halide; ring size; type of product}]
Which of the following is/are correct statement regarding major product?
\bfig{210.6}{TOAS17_p1_chem_s1_q02_1.png}
\begin{qoptions}
\opt{A}{Intramolecular nucleophilic substitution reaction takes place during the reaction sequence}
\opt{B}{5 membered ring is formed as major product}
\opt{C}{Grignard reagents (CH$_3$MgBr) prefer to give nucleophilic addition on polar $\pi$-bond}
\opt{D}{Cyclic ether is formed as final major product}
\end{qoptions}
\end{bankq}
\begin{bankq}[id=TOAS18-P2-CHEM-S3-Q06, source={TOAS - Test 18 (30 Oct 2022)}, paper=2, qno=6, type=single, subtopic={Properties}, answer={6}, answer_src=key, difficulty=1, diff_parts={K1 S1 I0 C0 F0}, diff_note={concepts: acetylation of hydroxyl groups | steps: mass gain per acetyl group; count}]
An organic compound A(C$_{10}$H$_{18}$O$_8$) on treatment with excess of CH$_3$COCl gives a fully acetylated product whose molar mass is found to be 518 g/mol. How many hydroxyl functional groups are present in A?
\end{bankq}

\section*{TAAS papers}
\begin{bankq}[id=TAAS01-P1-CHEM-S2-Q06, source={TAAS - Test 01 (18 Jun 2023)}, paper=1, qno=6, type=multi, subtopic={Preparation}, answer={A,B,C,D}, answer_src=key, difficulty=5, diff_parts={K2 S2 I1 C0 F1}, diff_note={concepts: cleavage of an ether, coupling of halides | steps: identify the ether from the Wurtz products; chlorides B and C; alkene D; hydration and hydroboration products}, also_in={TNAS - Test 04 (18 Jun 2023) P1 Q12}]
\bfig{159.6}{TAAS01_p1_chem_s2_q06_1.png}
\begin{qoptions}
\opt{A}{$D$ is CH$_3{-}$CH$=$CH$_2$}
\opt{B}{$E$ is \bfigi{47.7}{TAAS01_p1_chem_s2_q06_B.png}}
\opt{C}{$F$ is \bfigi{23.7}{TAAS01_p1_chem_s2_q06_2.png}}
\opt{D}{$B$ is \bfigi{33.3}{TAAS01_p1_chem_s2_q06_D.png}}
\end{qoptions}
\end{bankq}
\begin{bankq}[id=TAAS02-P2-CHEM-S2-Q05, source={TAAS - Test 02 (30 Jul 2023)}, paper=2, qno=5, type=multi, subtopic={Properties}, answer={A,D}, answer_src=key, difficulty=3, diff_parts={K1 S1 I1 C0 F1}, diff_note={concepts: periodic acid cleavage of a polyol | steps: cleavage of each vicinal diol; products}]
$\begin{array}{l}\mathrm{CH_2-OH}\\|\\\mathrm{CH-OH}\\|\\\mathrm{CH-OH}\\|\\\mathrm{CH_2-OH}\end{array}\ \xrightarrow[\text{excess}]{\mathrm{HIO_4}}$

which of the following can be formed by above reaction?
\begin{qoptions}
\opt{A}{$\mathrm{H-\overset{O}{\overset{\|}{C}}-H}$}
\opt{B}{CH$_3-$OH}
\opt{C}{CO$_2$}
\opt{D}{$\mathrm{H-\overset{O}{\overset{\|}{C}}-OH}$}
\end{qoptions}
\end{bankq}
\begin{bankq}[id=TAAS02-P2-CHEM-S2-Q06, source={TAAS - Test 02 (30 Jul 2023)}, paper=2, qno=6, type=multi, subtopic={Properties}, answer={A,B,C,D}, answer_src=key, difficulty=3, diff_parts={K1 S1 I1 C0 F1}, diff_note={concepts: cleavage of ethers | steps: mechanism for each ether with acid; match}]
Which of the following reaction are correctly matched?
\begin{qoptions}
\opt{A}{$\begin{array}{c}\mathrm{CH_3}\\|\\\mathrm{CH_3-C-O-CH_3}\\|\\\mathrm{CH_3}\end{array}\xrightarrow{\mathrm{HI}}\mathrm{S_N1}$}
\opt{B}{$\begin{array}{c}\mathrm{CH_3-CH-O-CH_3}\\|\\\mathrm{CH_3}\end{array}\xrightarrow{\mathrm{H_2O/H^+}}\mathrm{S_N1}$}
\opt{C}{$\begin{array}{c}\mathrm{CH_3-CH-O-CH_3}\\|\\\mathrm{CH_3}\end{array}\xrightarrow{\mathrm{HI}}\mathrm{S_N2}$}
\opt{D}{$\mathrm{CH_3{-}O{-}CH_2-CH_3}\xrightarrow{\mathrm{HBr}}\mathrm{S_N2}$}
\end{qoptions}
\end{bankq}

\section*{TMAS papers}
\begin{bankq}[id=TMAS01-P2-CHEM-S3-Q01, source={TMAS - Test 01 (09 Jul 2023)}, paper=2, qno=1, type=numerical, subtopic={Properties}, answer={7}, answer_src=key, difficulty=5, diff_parts={K1 S2 I1 C1 F1}, diff_note={concepts: oxidising agents | steps: nature of the conversion; reagents that work; count}]
How many reagents can be used for the given conversion.
\bfig{95.6}{TMAS01_p2_chem_s3_q01_1.png}
(1) Cu, 300$^\circ$C \quad (2) $\mathrm{KMnO_4/H^\oplus}$ \quad (3) PCC/$\mathrm{CH_2Cl_2}$ \quad (4) $\mathrm{K_2Cr_2O_7/H^\oplus}$\\
(5) $\mathrm{MnO_2}$ \quad (6) $\mathrm{Al(OCMe_3)_3}$/Acetone\\
(7) NBS \quad (8) (i) TsCl (ii) DMSO (iii) $\mathrm{NaHCO_3}$\\
(9) (i) DIBAL-H (ii) $\mathrm{H_2O}$
\end{bankq}
\begin{bankq}[id=TMAS02-P1-CHEM-S2-Q07, source={TMAS - Test 02 (06 Aug 2023)}, paper=1, qno=7, type=multi, subtopic={Properties}, answer={A,C}, answer_src=key, difficulty=6, diff_parts={K2 S2 I1 C1 F1}, diff_note={concepts: reaction sequence, alcohols | steps: products; acidity; Lucas test; unsaturation}]
$\mathrm{Ph-CHO}\xrightarrow[\mathrm{(ii)\,H^+/H_2O}]{\mathrm{(i)\,CH_3MgBr}}\mathrm{(A)}\xrightarrow{\mathrm{PCC}}\mathrm{(B)}$ \bfigi{131.1}{TMAS02_p1_chem_s2_q07_1.png}\par
Select the correct statement
\begin{qoptions}
\opt{A}{Acidic strength of $E$ is more than $D$}
\opt{B}{Both $D$ and $E$ gives immediate turbidity with Lucas reagent.}
\opt{C}{DOU of $D$ is $4$}
\opt{D}{Product $E$ is ethanol}
\end{qoptions}
\end{bankq}
\begin{bankq}[id=TMAS02-P2-CHEM-S1-Q06, source={TMAS - Test 02 (06 Aug 2023)}, paper=2, qno=6, type=multi, subtopic={Properties}, answer={A,C,D}, answer_src=solved, working={Periodate cleavage of the tetraol gives HCHO and HCOOH; excess periodate also oxidises HCOOH to $\mathrm{CO_2}$}, difficulty=5, diff_parts={K1 S2 I1 C1 F1}, diff_note={concepts: periodate cleavage | steps: products; further oxidation; check options}]
\bfig{65.6}{TMAS02_p2_chem_s1_q06_1.png}
which of the following can be formed by above reaction?
\begin{qoptions}
\opt{A}{\bfigi{38.5}{TMAS02_p2_chem_s1_q06_A.png}}
\opt{B}{$\mathrm{CH_3-OH}$}
\opt{C}{$\mathrm{CO_2}$}
\opt{D}{\bfigi{45.5}{TMAS02_p2_chem_s1_q06_D.png}}
\end{qoptions}
\end{bankq}

\section*{TNAS papers}
\begin{bankpassage}[id=TNAS02-P2-CHEM-S3-G1-2, source={TNAS - Test 02 (07 May 2023)}, questions={1-2}]
Acid catalysed conversion of 1,2-diol or vicinal diol, into carbonyl compound known as pinacol-pinacolone rearrangement.
\bfig{151.1}{TNAS02_p2_chem_s3_para1_2_1.png}
Generally more electron donating group migrate during the mechanism.
\end{bankpassage}
\begin{bankq}[id=TNAS02-P2-CHEM-S3-Q01, source={TNAS - Test 02 (07 May 2023)}, paper=2, qno=1, type=numerical, subtopic={Properties}, answer={3.00}, answer_src=key, difficulty=6, diff_parts={K1 S2 I2 C1 F1}, diff_note={concepts: pinacol rearrangement | steps: migration; product; unsaturation}, group=TNAS02-P2-CHEM-S3-G1-2]
Calculate the degree of unsaturation in the final product of given reaction?
\bfig{75.7}{TNAS02_p2_chem_s3_q01_1.png}
\end{bankq}
\begin{bankq}[id=TNAS02-P2-CHEM-S3-Q02, source={TNAS - Test 02 (07 May 2023)}, paper=2, qno=2, type=numerical, subtopic={Properties}, answer={2.00}, answer_src=key, difficulty=6, diff_parts={K1 S2 I2 C1 F1}, diff_note={concepts: pinacol type product | steps: migration; product structure; lone pairs}, group=TNAS02-P2-CHEM-S3-G1-2]
Number of lone pair present in the major product of the given reaction
\bfig{95.1}{TNAS02_p2_chem_s3_q02_1.png}
\end{bankq}
\begin{bankpassage}[id=TNAS02-P2-CHEM-S3-G3-4, source={TNAS - Test 02 (07 May 2023)}, questions={3-4}]
The mechanism of acid catalysed dehydration is given below
\bfig{124.1}{TNAS02_p2_chem_s3_para3_4_1.png}
This reaction take place in presence of conc. $\mathrm{H_2SO_4}$. In this reaction carbocation is formed and rearrangement may be possible.
\end{bankpassage}
\begin{bankq}[id=TNAS02-P2-CHEM-S3-Q03, source={TNAS - Test 02 (07 May 2023)}, paper=2, qno=3, type=numerical, subtopic={Properties}, answer={3.00}, answer_src=key, difficulty=5, diff_parts={K1 S2 I1 C1 F1}, diff_note={concepts: dehydration of 2-butanol | steps: alkenes possible; stereoisomers; count}, group=TNAS02-P2-CHEM-S3-G3-4]
Calculate the number of total alkene formed by dehydration of 2-butanol?
\end{bankq}
\begin{bankq}[id=TNAS02-P2-CHEM-S3-Q04, source={TNAS - Test 02 (07 May 2023)}, paper=2, qno=4, type=numerical, subtopic={Properties}, answer={3.00}, answer_src=key, difficulty=6, diff_parts={K1 S2 I2 C1 F1}, diff_note={concepts: rearrangement of carbocations | steps: each step; shifts; count}, group=TNAS02-P2-CHEM-S3-G3-4]
How many 1,2-shifts of carbocation intermediate are involved during the course of following reaction?
\bfig{225.3}{TNAS02_p2_chem_s3_q04_1.png}
\end{bankq}
\begin{bankq}[id=TNAS03-P1-CHEM-S2-Q04, source={TNAS - Test 03 (28 May 2023)}, paper=1, qno=4, type=numerical, subtopic={Properties}, answer={4}, answer_src=solved, working={Immediate turbidity needs a tertiary, benzylic or allylic secondary alcohol: tert-butyl alcohol, diphenylmethanol, 1-phenylethanol and cycloheptatrienol}, difficulty=5, diff_parts={K1 S2 I1 C1 F1}, diff_note={concepts: Lucas test | steps: classify the alcohols; immediate turbidity; count}]
$X$ = number of following alcohols give immidiate turbidity with Lucas reagent.\\
Find value of $X$
\bfig{124.9}{TNAS03_p1_chem_s2_q04_1.png}
(iii) $\mathrm{CH_3{-}OH}$ \quad (iv) $\mathrm{CH_3{-}CH_2{-}OH}$
\bfig{118.3}{TNAS03_p1_chem_s2_q04_2.png}
\bfig{98.0}{TNAS03_p1_chem_s2_q04_3.png}
\end{bankq}
\begin{bankq}[id=TNAS03-P2-CHEM-S2-Q07, source={TNAS - Test 03 (28 May 2023)}, paper=2, qno=7, type=single, subtopic={Properties}, answer={C}, answer_src=key, difficulty=5, diff_parts={K1 S2 I2 C1 F0}, diff_note={concepts: reactions of ethers and alcohols with halogen acids | steps: mechanism; products; match}]
Match the list-
\begin{center}
{\footnotesize\setlength{\tabcolsep}{2pt}\begin{tabular}{|c|p{0.42\linewidth}|c|p{0.26\linewidth}|}\hline
\multicolumn{2}{|c|}{List-I (Reactions)} & \multicolumn{2}{c|}{List-II (About major product)}\\\hline
(P) & \bfigi{71.5}{TNAS03_p2_chem_s2_q07_1.png}$\xrightarrow{\mathrm{HBr}}$ & (1) & Atleast one of the product is alcohol\\\hline
(Q) & \bfigi{74.0}{TNAS03_p2_chem_s2_q07_2.png}$\xrightarrow{\mathrm{HI}}$ & (2) & Atleast one product is alkyl halide\\\hline
(R) & \bfigi{58.7}{TNAS03_p2_chem_s2_q07_3.png}$\xrightarrow[\text{(excess)}]{\mathrm{HI}}$ & (3) & No reaction\\\hline
(S) & \bfigi{64.9}{TNAS03_p2_chem_s2_q07_4.png}$\xrightarrow[\mathrm{H_2SO_4}]{\mathrm{dil.}}$ & (4) & More than two organic product is obtained\\\hline
 & & (5) & Reaction follows nucleophilic substitution\\\hline
\end{tabular}}
\end{center}
Choose correct options
\begin{qoptions}
\opt{A}{$\mathrm P\to3,4;Q\to4,5;R\to1;S\to1,2,3$}
\opt{B}{$\mathrm P\to2,5;Q\to3;R\to2,5;S\to1,3,5$}
\opt{C}{$\mathrm P\to2,5;Q\to3;R\to2,4,5;S\to1,4,5$}
\opt{D}{$\mathrm P\to2,5;Q\to1,2,5;R\to2,4,5;S\to1,4,5$}
\end{qoptions}
\end{bankq}
\begin{bankq}[id=TNAS04-P2-CHEM-S1-Q07, source={TNAS - Test 04 (18 Jun 2023)}, paper=2, qno=7, type=multi, subtopic={Properties}, answer={A,B}, answer_src=key, difficulty=4, diff_parts={K1 S2 I1 C0 F1}, diff_note={concepts: epoxide ring opening with an alkyne anion | steps: attack at the less hindered carbon; alkoxide product; methylation}]
\bfig{126.8}{TNAS04_p2_chem_s1_q07_1.png}
Which of the following are correct?
\begin{qoptions}
\opt{A}{(P) is \bfigi{108.2}{TNAS04_p2_chem_s1_q07_2.png}}
\opt{B}{(Q) is \bfigi{108.3}{TNAS04_p2_chem_s1_q07_3.png}}
\opt{C}{(P) is \bfigi{108.3}{TNAS04_p2_chem_s1_q07_4.png}}
\opt{D}{(Q) is \bfigi{108.5}{TNAS04_p2_chem_s1_q07_5.png}}
\end{qoptions}
\end{bankq}
\begin{bankq}[id=TNAS05-P2-CHEM-S3-Q05, source={TNAS - Test 05 (09 Jul 2023)}, paper=2, qno=5, type=numerical, subtopic={Preparation}, answer={3}, answer_src=key, difficulty=5, diff_parts={K1 S2 I1 C1 F1}, diff_note={concepts: selective oxidation reagents | steps: nature of the conversion; reagents that stop at the aldehyde; reagents that cannot; count}]
Identify total number of reagents used for following conversion?
\bfig{104.1}{TNAS05_p2_chem_s3_q05_1.png}
(i) $\mathrm{KMnO_4}$ \quad (ii) $\mathrm{K_2Cr_2O_7}$ \quad (iii) $\mathrm{CrO_3.H_2O}$ \quad (iv) $\mathrm{CrO_3.Py.HCl/CCl_4}$ \quad (v) $\mathrm{SeO_2}$ \quad (vi) $\mathrm{MnO_2}$ \quad (vii) Tollen's reagent \quad (viii) Cu/300$^\circ$C
\end{bankq}
\begin{bankq}[id=TNAS07-P1-CHEM-S1-Q15, source={TNAS - Test 07 (20 Aug 2023)}, paper=1, qno=15, type=single, subtopic={Properties}, answer={B}, answer_src=key, difficulty=4, diff_parts={K1 S2 I2 C0 F0}, diff_note={concepts: multistep reaction product | steps: reaction at each step; intermediate; final product}]
\bfig{155.6}{TNAS07_p1_chem_s1_q15_1.png}
Final product is :
\begin{qoptions}
\opt{A}{\bfigi{21.1}{TNAS07_p1_chem_s1_q15_2.png}}
\opt{B}{\bfigi{25.8}{TNAS07_p1_chem_s1_q15_3.png}}
\opt{C}{\bfigi{25.8}{TNAS07_p1_chem_s1_q15_4.png}}
\opt{D}{\bfigi{23.3}{TNAS07_p1_chem_s1_q15_5.png}}
\end{qoptions}
\end{bankq}
\begin{bankq}[id=TNAS07-P1-CHEM-S2-Q09, source={TNAS - Test 07 (20 Aug 2023)}, paper=1, qno=9, type=single, subtopic={Preparation}, answer={1}, answer_src=key, difficulty=6, diff_parts={K1 S3 I2 C1 F0}, diff_note={concepts: Kolbe electrolysis of benzoate and the structure of biphenyl | steps: oxidation of the side chain; salt formation; Kolbe coupling; count sigma bonds; count pi electrons}]
\bfigi{66.4}{TNAS07_p1_chem_s2_q09_1.png}\ $X$\ \bfigi{69.4}{TNAS07_p1_chem_s2_q09_2.png}\par
$Y$(Aromatic hydrocarbon)\par
In `$Y$' total number of sigma bond formed by overlapping of ($sp^2-sp^2$) orbital(s) $=$ ($P$)\par
Total number of $\pi e^-$ in `$Y$' $=$ ($Q$)\par
value of $|P-Q|$ is
\end{bankq}
\end{document}
